\section{무엇을 고쳤나}

\paragraph{주장 1.} 「최적 base」는 자료의 사실이 아니라 격자의 사실이었다 --- 격자를 8 칸으로 늘리면 판정 자의 argmax 가 450 이고 격자 오른쪽 끝인가는 거짓 이며, argmax 와 2·SE_clu 로 안 갈리는 칸이 7 개다

\paragraph{주장 2.} 「무엇이 사는가」는 자의 사실이다 --- 자 3 × 성분 8 = 24 칸을 전량 재면 자에 따라 2·SE 판정이 갈리는 성분이 4 개이고, 상호작용은 판정 자에서 t_clu 1.467108 로 안 살고 챔피언 자에서 2.226044 로 산다

\paragraph{주장 3.} 변이체의 「공허」를 손 라벨로 판정하면 검정력이 0 이다 --- 앞 노트가 「구성상 거짓 0」이라 신고한 여섯을 그 노트 자신의 설정 격자에서 돌리면 6 이고, 신고값은 0 였다

\paragraph{주장 4.} 앞 노트가 손으로 센 「990 의 일곱 중 여섯」을 기계로 재면 자에 따라 7 과 6 로 갈린다 --- 「구성상」의 정의가 자를 가른다

\paragraph{주장 5.} 면제를 만든 사이클이 면제 없는 판을 안 내면 그 통과는 면제로 산 것이다 --- 앞 노트가 등기해 뺀 도장 셋이 앞 세 사이클에서 신판을 떨어뜨린 바로 그 셋인가: 참

\paragraph{주장 6.} 이 사이클의 배선 7 개는 전부 「설정에 따라 갈리는」 변이체를 가졌고 손 라벨이 0 이다 (검정력 있는 검사 7 · 구성상 거짓 0)


\section{한계}

이 스텝은 \textbf{컴파일하지 않았고 보내지 않았다}. 그 사실을 원장과 이 파일에 적는다 --- 「안 했다」는 「없다」가 아니다.
